% Part: methods
% Chapter: induction
% Section: relations

\documentclass[../../../include/open-logic-section]{subfiles}

\begin{document}

\olfileid{mth}{ind}{rel}

\olsection{رابطه‌ها و تابع‌ها}

وقتی مجموعه‌ای از اشیا، مانند اعداد طبیعی یا ترم‌های خوش‌ساخت، را
به‌طور استقرایی تعریف کرده باشیم، می‌توانیم \emph{رابطه‌هایی بر} این
اشیا را نیز با استقرا تعریف کنیم. برای نمونه، این ایده را در نظر
بگیرید: ترم خوش‌ساخت~$t_1$ زیرترم ترم خوش‌ساخت~$t_2$ است اگر به‌صورت
بخشی از آن ظاهر شود. نمادی برای این رابطه به کار می‌بریم:
$t_1 \sqsubseteq t_2$. البته هر ترم خوش‌ساخت زیرترم خودش است:
$t \sqsubseteq t$. می‌توانیم این رابطه را به‌طور استقرایی چنین تعریف
کنیم:

\begin{defn}
  رابطهٔ زیرترم‌بودن ترم خوش‌ساخت~$t_1$ از~$t_2$، یعنی
  $t_1 \sqsubseteq t_2$، با استقرا بر~$t_2$ چنین تعریف می‌شود:
  \begin{enumerate}
  \item اگر $t_2$ یک حرف باشد، آنگاه $t_1 \sqsubseteq t_2$ اگر و تنها
    اگر $t_1 = t_2$.
  \item اگر $t_2$ برابر $[s_1 \circ s_2]$ باشد، آنگاه
    $t_1 \sqsubseteq t_2$ اگر و تنها اگر $t_1 = t_2$، یا
    $t_1 \sqsubseteq s_1$، یا $t_1 \sqsubseteq s_2$.
  \end{enumerate}
\end{defn}

برای نمونه، این تعریف به ما می‌گوید
$\mathrm{a} \sqsubseteq [\mathrm{b} \circ \mathrm{a}]$. زیرا بند (۲)
می‌گوید $\mathrm{a} \sqsubseteq [\mathrm{b} \circ \mathrm{a}]$ اگر و
تنها اگر $\mathrm{a} = [\mathrm{b} \circ \mathrm{a}]$، یا
$\mathrm{a} \sqsubseteq \mathrm b$، یا
$\mathrm{a} \sqsubseteq \mathrm{a}$. دو مورد نخست کاذب‌اند:
$\mathrm{a}$ آشکارا با $[\mathrm{b} \circ \mathrm{a}]$ یکسان نیست، و
بنا بر (۱)، $\mathrm{a} \sqsubseteq \mathrm{b}$ اگر و تنها اگر
$\mathrm{a} = \mathrm{b}$، که این نیز کاذب است. اما باز هم بنا بر (۱)،
$\mathrm{a} \sqsubseteq \mathrm{a}$ اگر و تنها اگر
$\mathrm{a} = \mathrm{a}$، که صادق است.

باید توجه کرد که موفقیت این تعریف به واقعیتی وابسته است که هنوز اثبات
نکرده‌ایم: هر ترم خوش‌ساخت~$t$ یا خود یک حرف است، یا ترم‌های خوش‌ساخت
\emph{به‌طور یکتا تعیین‌شدهٔ} $s_1$ و $s_2$ وجود دارند به‌طوری‌که
$t = [s_1 \circ s_2]$. «به‌طور یکتا تعیین‌شده» اینجا یعنی اگر
$t = [s_1 \circ s_2]$، آنگاه $t$ نمی‌تواند \emph{همچنین} برابر
$[r_1 \circ r_2]$ باشد درحالی‌که $s_1 \neq r_1$ یا $s_2 \neq r_2$.
اگر چنین چیزی ممکن بود، بند~(۲) می‌توانست با خودش ناسازگار شود: با
خواندن $t_2$ به‌صورت $[s_1 \circ s_2]$ شاید به
$t_1 \sqsubseteq t_2$ برسیم، اما با خواندن $t_2$ به‌صورت
$[r_1 \circ r_2]$ شاید به این نتیجه برسیم که $t_1 \sqsubseteq t_2$
برقرار نیست. پیش از
اثبات ناممکن‌بودن این وضع، مثالی را ببینیم که در آن این اتفاق
\emph{ممکن} است.

\begin{defn}
  \emph{ترم‌های بی‌کروشه} را به‌طور استقرایی چنین تعریف می‌کنیم:
  \begin{enumerate}
  \item هر حرف یک ترم بی‌کروشه است.
  \item اگر $s_1$ و $s_2$ ترم‌هایی بی‌کروشه باشند، آنگاه
    $s_1 \circ s_2$ یک ترم بی‌کروشه است.
  \item هیچ چیز دیگری ترم بی‌کروشه نیست.
  \end{enumerate}
\end{defn}

برای مثال، $\mathrm{a}$، $\mathrm{b} \circ \mathrm{d}$ و
$\mathrm{b} \circ \mathrm{a} \circ \mathrm{b}$ ترم‌هایی بی‌کروشه‌اند.
اگر «زیرترم» را برای ترم‌های بی‌کروشه همانند بالا تعریف می‌کردیم، بند
دوم چنین می‌شد:
\begin{center}
  اگر $t_2 = s_1 \circ s_2$، آنگاه $t_1 \sqsubseteq t_2$ اگر و تنها
  اگر $t_1 = t_2$، یا $t_1 \sqsubseteq s_1$، یا
  $t_1 \sqsubseteq s_2$.
\end{center}

اکنون $\mathrm{b} \circ \mathrm{a} \circ \mathrm{b}$ به شکل
$s_1 \circ s_2$ است که در آن
\begin{align*}
  s_1 & = \mathrm{b} \text{ و } & s_2 & = \mathrm{a} \circ \mathrm{b}.
\intertext{همچنین به شکل $r_1 \circ r_2$ است که در آن}
  r_1 & = \mathrm{b} \circ \mathrm{a} \text{ و } & r_2 &= \mathrm{b}.
\end{align*}
آیا $\mathrm{a} \circ \mathrm{b}$ زیرترم
$\mathrm{b} \circ \mathrm{a} \circ \mathrm{b}$ است؟ اگر خوانش نخست
را مبنا بگیریم پاسخ مثبت، و اگر خوانش دوم را مبنا بگیریم پاسخ منفی است.

خاصیتی که می‌گوید شیوهٔ ساخته‌شدن یک ترم خوش‌ساخت از ترم‌های خوش‌ساخت
دیگر یکتاست، \emph{خوانش یکتا} نام دارد. چون تعریف استقرایی رابطه‌ها
برای چنین اشیای به‌طور استقرایی تعریف‌شده‌ای مهم است، باید اثبات کنیم
این خاصیت برقرار است.

\begin{prop}
  فرض کنید $t$ ترمی خوش‌ساخت است. آنگاه یا $t$ خود یک حرف است، یا
  ترم‌های خوش‌ساخت یکتای $s_1$ و $s_2$ وجود دارند به‌طوری‌که
  $t = [s_1 \circ s_2]$.
\end{prop}

\begin{proof}
  اگر $t$ خود یک حرف باشد، شرط برقرار است. پس فرض کنید $t$ خود یک حرف
  نیست. از تعریف استقرایی می‌دانیم که در این صورت $t$ باید برای برخی
  ترم‌های خوش‌ساخت $s_1$ و~$s_2$ به شکل $[s_1 \circ s_2]$ باشد. فقط
  باید نشان دهیم این دو به‌طور یکتا تعیین شده‌اند؛ یعنی اگر
  $t = [r_1 \circ r_2]$، آنگاه $s_1 = r_1$ و $s_2 = r_2$.

  پس فرض کنید برای ترم‌های خوش‌ساخت $s_1$، $s_2$، $r_1$ و $r_2$، هم
  $t = [s_1 \circ s_2]$ و هم $t = [r_1 \circ r_2]$. باید نشان دهیم
  $s_1 = r_1$ و $s_2 = r_2$. نخست، $s_1$ و $r_1$ باید یکسان باشند،
  زیرا در غیر این صورت یکی پاره‌رشتهٔ آغازی سره‌ای از دیگری است. اما
  بنا بر \olref[sti]{prop:initial}، وقتی $s_1$ و~$r_1$ هر دو ترم
  خوش‌ساخت باشند این امر ناممکن است. اگر $s_1 = r_1$، آنگاه به‌روشنی
  $s_2 = r_2$ نیز برقرار است.
\end{proof}

همچنین می‌توانیم تابع‌ها را به‌طور استقرایی تعریف کنیم. برای نمونه،
تابع~$f$ را که هر ترم خوش‌ساخت را به بیشترین عمق $[\dots]$های تودرتوی
آن می‌فرستد چنین تعریف می‌کنیم:

\begin{defn}
  \ollabel{defn:depth} \emph{عمق} یک ترم خوش‌ساخت، یعنی $f(t)$، به‌طور
  استقرایی چنین تعریف می‌شود:
  \[
  f(t) = \begin{cases}
    0 & \text{ اگر $t$ یک حرف باشد}\\
    \max(f(s_1), f(s_2)) + 1 & \text{ اگر $t = [s_1 \circ s_2]$.}
  \end{cases}
  \]
\end{defn}

برای نمونه،
\begin{align*}
  f([\mathrm{a} \circ \mathrm{b}])
    &= \max(f(\mathrm{a}),f(\mathrm{b})) + 1 \\
    &= \max(0, 0) + 1 = 1, \text{ و}\\
  f([[\mathrm{a} \circ \mathrm{b}] \circ \mathrm{c}])
    &= \max(f([\mathrm{a} \circ \mathrm{b}]), f(\mathrm{c})) + 1 \\
    &= \max(1,0) + 1 = 2.
\end{align*}

البته اینجا فرض می‌کنیم $s_1$ و $s_2$ ترم‌هایی خوش‌ساخت‌اند و از این
واقعیت استفاده می‌کنیم که هر ترم خوش‌ساخت یا یک حرف است یا به شکل
$[s_1 \circ s_2]$. باز هم مهم است که ترم فقط به یک شیوه می‌تواند چنین
شکلی داشته باشد. برای دیدن چرایی آن، دوباره ترم‌های بی‌کروشه‌ای را که
پیش‌تر تعریف کردیم در نظر بگیرید. «تعریف» متناظر چنین خواهد بود:
\[
  g(t) =
  \begin{cases}
    0 & \text{ اگر $t$ یک حرف باشد}\\
    \max(g(s_1), g(s_2)) + 1 & \text{ اگر $t = s_1 \circ s_2$.}
  \end{cases}
\]
اکنون ترم بی‌کروشهٔ
$\mathrm{a} \circ \mathrm{b} \circ \mathrm{c} \circ \mathrm{d}$ را
در نظر بگیرید. می‌توان آن را به بیش از یک شیوه خواند؛ برای مثال، به
صورت $s_1 \circ s_2$ که در آن
\begin{align*}
  s_1 & = \mathrm{a} \text{ و } &
  s_2 & = \mathrm{b} \circ \mathrm{c} \circ \mathrm{d},
\intertext{یا به‌صورت $r_1 \circ r_2$ که در آن}
  r_1 & = \mathrm{a} \circ \mathrm b \text{ و } &
  r_2 &= \mathrm{c} \circ \mathrm{d}.
\end{align*}
محاسبهٔ $g$ بر پایهٔ خوانش نخست می‌دهد:
\begin{align*}
  g(s_1 \circ s_2)
    &= \max(g(\mathrm{a}), g(\mathrm{b} \circ \mathrm{c} \circ \mathrm{d})) + 1\\
    &= \max(0,2) + 1 = 3
  \intertext{درحالی‌که بر پایهٔ خوانش دیگر داریم}
  g(r_1 \circ r_2)
    &= \max(g(\mathrm{a} \circ \mathrm{b}), g(\mathrm{c} \circ \mathrm{d})) + 1\\
    &= \max(1,1) + 1 = 2.
\end{align*}
اما یک تابع باید همیشه مقداری یکتا بدهد؛ پس «تعریف» ما از~$g$ اصلاً
تابعی را تعریف نمی‌کند.

\begin{prob}
  تعریفی استقرایی از تابع $l$ ارائه کنید، که در آن $l(t)$ تعداد نمادها
  در ترم خوش‌ساخت~$t$ است.
\end{prob}

\begin{prob}
  با استقرای ساختاری بر ترم‌های خوش‌ساخت $t$ اثبات کنید
  $f(t) < l(t)$، که در آن $l(t)$ تعداد نمادهای~$t$ است و $f(t)$ عمق
  $t$ است، چنان‌که در \olref[mth][ind][rel]{defn:depth} تعریف شد.
\end{prob}

\end{document}
